Allora, condividiamo l'ultima presentazione. Questa è molto veloce e poi veniamo a argomento. Questa è la parte finale. Abbiamo parlato del trasporto del segnale, della modulistica elettronica che ci può servire. Passiamo alla acquisizione dei dati. In parte abbiamo visto, all'inizio dell'anno, quel programmino per acquisire dati da sensori. Quello è un sistema di acquisizione dati. Abbiamo dei segnali che noi acquisivamo con uno opportuno software. Abbiamo utilizzato dei software in laboratorio per acquisire i dati provenienti da una ad una DC. Qualcosina abbiamo visto. Facciamo un discorso un po' più generico. Allora, sull'acquisizione di dati, quindi sui sistemi di acquisizione dati, chiaramente tutto parte dall'utilizzo di sensori che vanno a misurare grandezze fisiche. Un po' possiamo immaginare come quello che fa l'esempio di nostro occhio, che traduce dei fotoni, degli impulsi luminosi anche qui, degli impulsi elettri, che sono inseriti giustanzialmente. Il problema è il seguente. Se io ho una grandezza fisica che sta cambiando nel tempo e che io sto misurando con un sensore, questo sensore chiaramente fornirà un'informazione di tipo elettrico, quindi ad esempio una tensione, una carica, comunque qualcosa che io posso poi andare a misurare, una corrente, che varierà nel tempo. Quindi immaginiamo ad esempio un sensore di intensità luminosa, sto misurando l'intensità luminosa di un ambiente, e questa intensità luminosa con questo sensore cambia nel tempo. E io voglio vedere come varie nel tempo, è quello che avete fatto con Arduino sostanzialmente, avete deciso di utilizzare, ad esempio, il sensore di intensità luminosa e di campionare, con una certa frequenza di campionamento, il signore che fuoriesce via questo sensore. Capite che il campionamento viene effettuato con una certa frequenza di campionamento. Quindi se ad esempio siete interessati a fenomeni che variano lentamente nel tempo, immaginate di voler valutare come varia l'intensità luminosa in funzione nell'arco della giornata, quindi variazioni giorno-notte. È chiaro che non ha senso in questo caso a cui si è redato ogni secondo, perché tra un secondo e l'altro la variazione intensità luminosa è praticamente trascurabile. In quel caso non arriva bene anche misurarlo una volta ogni 10 minuti. Ma se invece siete interessati a percepire altri tipi di variazioni più brevi, come possono essere quelle legate, ad esempio la frequenza direte, i 50 Hz della frequenza direte che si riflette, ad esempio, nelle lampade del soffitto, o in una lampada qualsiasi in crescenza. Allora, in quel caso, capite che la frequenza di campionamento che dovete adoperare deve essere più elevata, proprio per percepire queste variazioni così frequenti. Quindi quando vi trovate a valutare, a acquisire i dati, la prima domanda che vi dovete porre è quale è la frequenza che io devo adoperare per rappresentare in maniera opportuna il fenomeno fisico che voglio studiare. Quello che stiamo facendo, sostanzialmente, è passare da una variazione continua della grandezza fisica a una variazione, in questo caso, discreta. Quindi io rappresento la mia grandezza fisica che varia con continuità attraverso un certo numero di punti. Più elevato questo numero di punti, più la rappresentazione sarà fedele. Questo è un processo che prende il nome di digitalizzazione. Quindi la digitalizzazione consiste nel campionare una determinata grandezza fisica nel tempo. Poter appunto rappresentare in funzione del tempo. Quindi memorizzare i dati e poter realizzare anche delle analisi successive di quello che è stato acquisito. Perché è opportuno, è utile, andare passando dal mondo analogico al mondo digitale. Un vantaggio è che abbiamo i dati che possono essere processati attraverso opportuni software di analisi. In più le quantità numeriche, quindi numeri, non sono più sensibili a rumore, cross talk, interferenze, cosa che invece caratterizza la trattazione di un segnale netto. Inoltre si possono memorizzare informazioni in modo permanente. Tuttavia, lo svantaggio è che chiaramente perdiamo fedeltà nel rappresentare il segnale, possiamo dire. Vedete, rispetto a un continuo, chiaramente noi ci limitiamo sempre un numero limitato di punti. E, oltre alla quantizzazione orizzontale, quegli a quanti punti io riesco ad accedere, ragionevolmente, abbiamo un problema anche di quantizzazione verticale, cioè questa misura che io effettuo, come viene effettuata. Ovviamente con un ADC devo valutare l'ampiezza di un segnale in un determinato istante. E l'ADC è una sua risoluzione. Quindi abbiamo un'indeterminazione, diciamo, orizzontale dovuta alla frequenza di campionamento e una verticale dovuta alla precisione con cui io vado a effettuare la misura. La digitalizzazione quindi può anche modificare in modo sostanziale la forma del segnale. Vediamo qui un esempio, un segnale continuo viene preso per punti e quindi questa è la sua rappresentazione digitale. Quindi ci dobbiamo accontentare, abbiamo dei vantaggi, ma ci dobbiamo accontentare del fatto che il segnale possa essere rappresentato più o meno bene. Che vi dicevo in base al fenomeno che poi uno vuole andare a studiare, si decide che frequenza ad operare. 

Ad esempio, immaginate di dover ascoltare della musica, di dover campionare, quindi la musica è trasformarla in un file. Chiaramente tutto dipende dalla frequenza con cui è andata a campionare la musica, quindi la qualità del suono che è andata a ascoltare, dipende dalla frequenza con cui è sempre campionata. Quindi a seconda del campo si decide di adoperare frequenze di campionamento diverse. Queste sono delle informazioni che già abbiamo visto, riguardanti il ADC, perché vi dicevo questo avvia attraverso comunque sia degli ADC. Abbiamo avuto queste cose. A questo punto, quando abbiamo a che fare con segnali impulsivi, capite che questi segnali normalmente hanno delle durate molto brevi. Quindi non è facile andare a digitalizzare un segnale che segnale che segnale che impulso, proprio niente da un rivelatore. Normalmente noi ci limitiamo a utilizzare l'ADC per misurare solo l'ampiezza massima. Tuttavia ad oggi esistono anche dei digitalizzatori con frequenze di campionamento sufficientemente levate per poter campionare tutto il segnale nel tempo. Ma queste sono applicazioni molto particolari. Normalmente non si è interessati ad acquisire tutta la forma d'onda. Un aspetto invece importante riguarda il fatto che oltre ad andare ad acquisire una sola variabile una sola grandezza fisica si può essere interessati in un dato evento a misurare simultaneamente più grandezza fisica. Quindi l'evento che è andato ad acquisire in realtà corrisponde a più informazioni simultaneamente. Questo è un esempio di acquisizione dati multi parametrica. Noi non ne vedremo dal vivo, il laboratorio, però è importante capire il concetto, la differenza. Quindi quello che fate voi abitualmente è acquisire una sola grandezza fisica per volta. Può essere ad esempio l'energia depositata in un rivelatore. Potrebbe ad esempio che un po' quello che avete fatto con questi moduli ADC producendo lo spetto. in avanti, potrebbe essere interessante andare ad acquisire più informazioni. Ad esempio se veste due rivelatori attraversate dalla stessa particella, potrebbe essere interessata ad esempio a misurare simultaneamente l'energia depositata nel primo rivelatore. L'energia depositata nel secondo rivelatore è il tempo di distanza di tempo tra i due segnali. Quindi un singolo evento in realtà contiene più informazioni. Energia persa nel primo, energia persa nel secondo, differenza di tempo. Magari quello che si fa è acquisire imponendo un trigger dato dalla coincidenza in maniera tale da essere sicuri di avere informazioni sensate in questo tipo di file di dati. Ma la cosa importante è sapere gestire l'analisi di questi dati. Capite che diventa tutto molto più complesso perché fino ad adesso la vada che fare con una sola variabile, adesso il problema diventa multidimensionale perché si potreste andare a guardare simbolarmente lo spettro in energia nel primo rivelatore, lo spettro in energia nel secondo oppure fare la somma dei due e guardare l'energia complessiva. Ma potrebbe essere invece interessante guardare come l'energia persa nel primo rivelatore dipende da tì, quindi realizzare un grafico bidimensionale o ancora grafici tridimensionali o guardare una variabile imponendo un talido su un'altra. Quindi capite che diventa un tipo di analisi molto diverso rispetto a quella cui siamo abituati normalmente. Come si fa l'analisi multiparametrica la posizione dati multiparametrica veramente dovete avere a disposizione dei sistemi che vi permettono di farlo. Quello che avete utilizzato per il laboratorio era un sistema che prevedeva un solo canale sostanzialmente quindi una sola variabile. Per fare questo ci viene in aiuto, c'è un comodo in aiuto dei sistemi un po' più complessi che prendere il nome torno mattine indietro ad esempio di sistemi VMI dove vedete nuovamente abbiamo sempre dei crate dove si vanno a inserire dei moduli e questa volta i crate sono collegati al PC e permettono un'acquisizione di più date simultaneamente oppure uno standard un po' più pecchiotto e lo standard camac che quello che vedete qui, ma sostanzialmente il principio lo stesso ci utilizzare dei moduli che ci permettono di avere ad esempio più a 10 più più TDC simultaneamente e poter trasmettere i dati al computer attraverso un bus la differenza tra i due questo è un po' più vecchio questo è un po' più recente che qui invece possiamo avere delle CPU a bordo quindi c'è anche dell'intelligenza a bordo e direttura i dati potrebbero rimanere proprio sul WMI, fare delle operazioni estralizzati e poi solamente trasferiti alla fine ora vi dicevo questo nello specifico non lo vedremo in laboratorio però è utile aver presente soprattutto per quanto riguarda l'analisi dei dati che a volte in fisica vi troverete a che fare con dati che richiederanno un'analisi richiederanno un'analisi multi parametrica ora il route questo si può fare con degli opportuni oggetti che prende il nome di 3 sicuramente in altri programmi come può essere matlab o origine o non li conosco tutti questo si farà in maniera alternativa anche con Python si possono fare e quindi quello che faremo successivamente sarà appunto provare a introdurre un software di analysis in questo caso route per affrontare questo tipo di problematica quindi partire da analisi semplici come possono essere istogrammi, grafici, fit faremo anche i numeri pseudo casuali quindi simulazioni e ci arriviamo vediamo un po' anche come evolvere nel tempo faremo anche un piccolo esempio di analisi dati multi parametrica su dei dati in questo caso simulati perché non abbiamo dei dati reali quindi io vi dirò quando inizieremo la lezione prima quando inizieremo corrutta